\documentclass[11pt]{article}
\usepackage{amsmath,amsfonts}
\usepackage{graphicx}
%\usepackage{xcolor}
%\definecolor{gray}{rgb}{.75,.75,.75}
%\usepackage[landscape]{geometry}
%\usepackage{hyperref}
%\usepackage[bookmarks=true, pdfborder={0 0 .5}, urlbordercolor=gray]{hyperref}

\renewcommand{\thefootnote}{\fnsymbol{footnote}}


\addtolength{\topmargin}{-.65in}
\addtolength{\textheight}{1.3in}
%\addtolength{\oddsidemargin}{-.8in}
%\addtolength{\textwidth}{1.4in}
\pagestyle{empty}

\newcommand{\ds}{\displaystyle}
\newcommand{\ts}{\textstyle}

\newcommand{\Z}{\mathbb{Z}}
\newcommand{\R}{\mathbb{R}}
\newcommand{\C}{\mathbb{C}}
\newcommand{\EE}{\mathcal{E}}
\newcommand{\tZ}{\tilde{\Z}}

\newcommand{\Sym}{\mathrm{Sym}}

\renewcommand{\circle}[1]{{\bigcirc\hspace{-.75em}#1}}

\begin{document}

\footnote[0]{Notes by Peter Magyar \ {\tt magyar@math.msu.edu}}
\vspace{-1em}

\centerline{\bf Math 133\hspace{1.5em} \hfill{\Large\bf Volume}\hfill Stewart \S5.2}
\vspace{1em}

\noindent
{\bf Geometry of integrals.} In this section, we will learn how to compute volumes 
using integrals defined by slice analysis. First, we recall from Calculus I how to compute the area
under a graph as the cumulative effect (the integral) of the height of the graph.
Given the region under $y=f(x)$ and above an interval $[a,b]$ on the $x$-axis,
we slice up the interval into $n$ increments (small parts) of width $\Delta x=\frac{b-a}{n}$,
with division points:
$$
a < a{+}\Delta x <a{+}2\Delta x<\cdots < a{+}n\Delta x=b,
$$
and we take sample points $x_1,\ldots, x_n$, one in each increment.
This slices the area into increments $\Delta A_1,\ldots, \Delta A_n$, 
each approximately a rectangle:
\\[1em]
\centerline{
\includegraphics[width=2.2in]
{5-2-Area.png}
}
\vspace{-1em}

\noindent
The area of an increment is approximately:
$$
\Delta A_i \ \approx\ \text{(height)$\times$(width)} \ =\ f(x_i)\,\Delta x\,,
$$
so the total area is the Riemann sum (see \S4.1 Pt 1):
$$
A\ \ =\ \sum_{i=1}^n \Delta A_i
\approx\ 
\sum_{i=1}^n f(x_i)\Delta x\,.
$$
\vspace{-.5em}

\noindent 
Taking the limit of very many thin slices gives the exact area as the integral:
$$
A 
\ =\ \lim_{n\to\infty} \sum_{i=1}^n f(x_i)\Delta x\ =\ \int_a^b\! f(x)\,dx\,.
$$
But this only gives approximate numerical answers to our question, 
unless we can determine the limit, which is very difficult directly (\S4.1 Pt 2). 

No really surprising ideas so far: 
for thousands of years, mathematicians used similar methods to compute areas, 
with very limited success.
The genius of Newton shines in three ideas which solve the problem, 
easily computing almost any area (\S4.3). 
The first idea is to make area into a function: $A(x)=\int_a^x\! f(t)\,dt$ is the area
over a variable interval $t\in [a,x]$. 
Second, since area is the cumulative effect of the height $y=f(x)$,
the rate of change of area is the height at the moving endpoint:
$A'(x) = f(x)$. 
This is the First Fundamental Theorem of Calculus: $A(x)$ is an antiderivative of $f(x)$. 

The third and clinching idea is the algebra of differentiation.
If we can reverse the derivative rules to find another antiderivative, 
a known function $F(x)$ with $F'(x)=f(x)$, 
then the Uniqueness Theorem (\S3.2) tells us 
that the two antiderivatives are the same, 
except for adding a constant:
$
A(x) = \int_a^x\! f(t)\,dt \ =\ F(x)+C
$.
Since $0=A(a)=F(a)+C$, we have $C=-F(a)$, and $A(x)=F(x)-F(a)$,  
This is the Second Fundamental Theorem:
$$
A = A(b) = \int_a^b\! f(x)\,dx \ =\ F(b)-F(a).
$$
That is, the area is the integral of a {\it rate of change} $f(x) = F'(x)$, which is
is equal to the {\it total change} in $F(x)$ over $x\in[a,b]$.\footnote{
The {\it indefinite integral} $\int f(x)\,dx = F(x)+C$ is 
the general antiderivative function, also called the primitive function.
The {\it definite integral} $\int_a^b\! f(x) = F(b)-F(a)$ is a number.
}

These three brilliant ideas, working perfectly together, make it easy to compute the areas of most shapes, provided we are skillful enough at reversing derivative rules to find indefinite integrals.
\\[.5em]
{\sc example:} Find the area of the leftmost region enclosed below the curve $y=x \cos(x^2)$ and above the line  $y=-x$.
\\[1em]
\centerline{
\includegraphics[width=1.3in]
{5-2-XCosX2.png}
}
\vspace{-.5em}

\noindent
The left corner of the region is at $x=0$.
To find the right corner, we set $x\cos(x^2)=-x$ with $x>0$, so $\cos(x^2)=-1$, 
and the smallest positive solution is $x^2=\pi$, $x=\sqrt{\pi}\approx 1.77$\,. 

To compute the area, we consider thin slices from the top curve to the bottom line, with height $x\cos(x^2)-(-x)=x\cos(x^2)+x$, and take a limit of Riemann sums to get:
$$
A
\ =\ 
\int_0^{\!\sqrt{\pi}}\!\! \left(x\cos(x^2)+x\right)dx
\ =\ \int_0^{\!\sqrt{\pi}}\!\!  x\cos(x^2)\,dx\ +\
\int_0^{\!\sqrt{\pi}}\!\!  x\,dx.
$$
We can evaluate the first term using the substitution $u=x^2$, $du=2x\,dx$, so
$\int\! x\cos(x^2)\,dx$ $=\frac12\int\cos(x^2)\, 2x\,dx=\frac12\int \cos(u)\,du=\frac12\sin(u)=\frac12 \sin(x^2)$.
$$
A 
\ =\
\left[\vphantom{\sum}\tfrac12\sin(x^2) + \tfrac12 x^2\right]_{x=0}^{x=\sqrt{\pi}}
\ = \
\tfrac12{\pi}\ \approx\ 1.57\,.
$$
This is the same as the area of a semi-circle of radius 1, which looks about right from the picture. Not a result you could get from elementary geometry!

\pagebreak

\noindent
{\bf Solids of revolution.} The integral is a very powerful tool to compute the total of many small parts, such as the thin slices which fill up an area. 
%To get an exact result, we take the limit of very many, very thin slices.
A similar method allows us to compute volumes. 
We start with a {\it solid of revolution}
which is formed by rotating a region in the plane around the $x$-axis, sweeping out a kind of barrel shape.
\\[1em]
\centerline{
\includegraphics[width=2.2in]
{5-2-Region.png}
\qquad \ \ 
\includegraphics[width=2.2in]
{5-2-Rotated-Region.png} 
}
\vspace{0em}

\noindent
The total volume is  the sum of $n$ disk slices at sample points $x_1,\ldots,x_n$ in $[a,b]$, each with radius $f(x_i)$ and thickness $\Delta x=\frac{b-a}{n}$.
\\[1em]
\centerline{
\includegraphics[width=2.4in]
{5-2-Rotated-Disk.png} 
}
\vspace{-.3em}

\noindent
The volume of each slice is approximately:
$$
\Delta V_i \ \approx\ \text{(circle area)$\times$(thickness)} 
\ =\ \pi f(x_i)^2\, \Delta x\,,
$$
and taking the limit as $n\to\infty$ gives the exact volume as an integral:
$$
V \ =\ \lim_{n\to\infty} \sum_{i=1}^n \pi f(x_i)^2\,\Delta x
\ =\ \int_a^b \pi f(x)^2\,dx\,.
$$
Once we express our answer as an integral, we no longer consider its geometric motivation: finding an antiderivative and determining the value is a purely algebraic problem.

This analysis should not be surprising: we expect our answer to converge to an integral as soon as we can slice up the problem into small increments.
\\[.5em]
{\sc example:} Find the volume of the trumpet solid obtained by rotating around the $x$-axis the region defined by:
$0\leq y\leq \frac1x$ and $1\leq x\leq 2$. 
This is the solid of revoltion of the curve $y=\frac1x$ over the interval $x\in[a,b]$. 
\\[1em]
\centerline{
\includegraphics[width=2in]
{5-2-Trumpet.png} 
}
\vspace{0em}

\noindent
Applying our formula: 
$$
V
\ =\ 
\int_1^2\pi\left(\tfrac1x\right)^2\,dx
\ =\ 
\int_1^2\pi x^{-2}\,dx
$$
$$
\ =\ 
\left[-\pi x^{-1}\right]_{x=1}^{x=2}
\ =\ 
(-\pi 2^{-1})-(-\pi 1^{-1})
\ = \ \tfrac{\pi}{2}\,.
$$
\\[-.5em]
{\sc example:} Find the volume of the solid obtained by rotating 
the region defined by
$x^2\leq y\leq 2$ and $x\geq 1$,
around the vertical axis $x=1$.
\\[1em]
\centerline{
\includegraphics[width=3in]
{5-2-Y-Rotation.png} 
}
\vspace{0em}

\noindent
The setup is different from the $x$-axis rotation, so
we must repeat our slice analysis. 
Since we rotate around the vertical line $x=1$, 
we should take horizontal slices positioned by their height $y$,
and we should write our region as lying next to the interval
$y\in [1,2]$ between the curves $x=1$ and $x=\sqrt y$.

Our slices are thin horizontal disks at sample heights 
$y_1,\ldots, y_n$ in the interval $y\in[1,2]$. 
The thickness of each disk is $\Delta y$.
The radius at height $y_i$ is the horizontal distance 
from the axis $x=1$ to the curve $x=\sqrt y$;
this distance is $\sqrt{y_i}-1$. The volume of each disk is 
$\Delta V_i\approx\pi (\sqrt{y_i}{-}1)^2\,\Delta y$, and:
$$\hspace{-.5in}
V
\ =\ 
\int_1^2\!\!\pi\!\left(\sqrt{y}{-}1\right)^2\,dy
\ =\ 
\int_1^2\!\!\pi(y{-}2\sqrt{y}{+}1)\,dy
\ =\ 
\left[\pi\!\left(\tfrac12 y^2{-}\tfrac43y^{3/2}{+}y\right)\right]_{y=1}^{y=2}
\ =\ \pi\! \left(\tfrac{23}{6}{-}\tfrac{8\sqrt2}3\right)\!.
$$
%$$
%\ =\ 
%\pi\!\left(\tfrac12 2^2{-}\tfrac43 2^{3/2}{+}2\right)
%- \pi\!\left(\tfrac12 1^2{-}\tfrac43 1^{3/2}{+}1\right)
%\ = \ 
%\pi\! \left(\tfrac{23}{6}-\tfrac{8\sqrt2}3\right)
%\ \approx\ 0.195
%$$

\pagebreak

\noindent
{\bf Solids with specified cross sections.} Consider a solid rather like the sail 
of the man-o-war jellyfish. It has a horizontal base outlined by the circle
$x^2+y^2= 1$ in the $xy$-plane, and its vertical cross section above each line $x=c$ 
is an isosceles right triangle, with height equal to half its base length.
\\[.5em]
\centerline{
\includegraphics[width=2.8in]
{5-2-Man-O-War.png} 
}
\vspace{-1em}

\noindent
The slices are defined over $x\in[-1,1]$, and each slice at $x=x_i$ has a base segment
reaching from $y=-\sqrt{1-x^2}$ to $y=\sqrt{1-x^2}$, a thickness $\Delta x$,
and an area $\frac12$(base)$\times$(height) = $\tfrac12\!\left(\!2\sqrt{1{-}x_i^2}\,\right)\!\!\left(\!\sqrt{1{-}x_i^2}\,\right)$.
Thus:
$$
V
\ =\ 
\int_{-1}^1\!\tfrac12\!\left(\!2\sqrt{1{-}x^2}\right)\!\!\left(\!\sqrt{1{-}x^2}\right)\,dx
\ =\ 
\int_{-1}^1\!\! 1{-}x^2\,dx
\ =\ 
\left[\vphantom{\sum}x {-} \tfrac13x^3\right]_{x=-1}^{x=1}
 =\ \tfrac43.
$$
\\[0em]
{\bf Method of slice analysis.} 
The goal is to compute the size or bulk $S$ of a geometric object: 
its length, area, volume, mass, energy, etc.
\begin{enumerate}
\item Cut the object into thin slices whose position is determined by
some variable $x\in[a,b]$, often slices perpendicular to the $x$-axis. 

\item Compute the approximate size of the slice with thickness $\Delta x$ 
near a particular value of $x$ as: $\Delta S \approx f(x)\,\Delta x$.
Summing up slices gives the total size, and the limit of thinner slices gives the integral:
$$
S  =\ \int_a^b\!\! f(x)\,dx\,.
$$
The size is the accumulation of the slice measurement $f(x)$: 
e.g.~area is accumulation of slice height, volume is accumulation of slice area.

\item Evaluate the integral $S=\int_a^b\!f(x)\,dx$ algebraically or numerically.

\end{enumerate}
\bigskip

\noindent
{\small
{\sc challenge problem:} Consider the solid of revolution around the $x$-axis of the curve $y=f(x)$
over the interval $x\in[a,b]$. Explain the following formula for the surface area of 
this solid: $$S \ =\ \int_a^b 2\pi f(x)\sqrt{1{+}(f'(x))^2\,}\, dx\,.$$
Hint: \
$\sqrt{(\Delta x)^2 + (\Delta y)^2\,} 
\ =\ \sqrt{1+(\frac{\Delta y}{\Delta x})^2\,}\ \Delta x 
\ \approx\ \sqrt{1+(\frac{dy}{dx})^2\, }\ \Delta x$.
}

\end{document}



\begin{itemize}

\item Volume of rotation ($x$-axis): \\
disk slices picture $\Longrightarrow$ sum formula
$\Longrightarrow$ $V=\int_a^b \pi f(x)^2\,dx$

\item example: Trumpet $0\leq y\leq 1/x$ for $1\leq x\leq 2$ around $x$-axis:\\
$V =\int_1^2 \pi\!(1/x)^2\,dx = \left[-x^{-1}\right]_{x=1}^{x=2} = (-2^{-1}) - (-1^{-1}) =\frac12$.

\item example: Hollow trumpet $1/x^2\leq y\leq 1/x$ and $1\leq x\leq 2$,
$x$-axis:\\
integral of washers = difference of solids of revolution
\\
$V =\int_1^2 \pi(1/x)^2-\pi(1/x^2)^2 \,dx = \left[-x^{-1}+\frac13 x^{-3}\right]_{x=1}^{x=2}  = (-2^{-1}) - (-1^{-1}) 
(-\frac13 2^{-3} - (-\frac13 1^{-3})=\frac5{24}$.

\item Other axes: $y$-axis $V=\int_c^d \pi g(y)^2\,dy$\\
$(y=c)$-axis: $V=\int_a^b \pi (f(x){-}c)^2\,dx$

\item Example: $x^2\leq y\leq 2$ $x\geq 1$ $y=1$-axis .\\
$y=x^2\Rightarrow x =\sqrt y$\\
$V = \int_1^2 \pi (\sqrt{y}-1)^2\,dy =\left[\frac12 y^2 -2(\frac23 y^{3/2}) + y\right]_{y=1}^{y=2} \approx 0.6  $

\item Figure defined by fixed base and similar vertical slices:\\
--Fin shape: circle base, triangular slices: picture with Paint \\
\indent $V=\int_{-1}^1 \frac12 (2\sqrt{1{-}x^2})(\sqrt{1{-}x^2})\,dx=\int_{-1}^1 1-x^2\,dx$
\\
--Sphere: circle base, semi-circle slices $\Longrightarrow$ 
\end{itemize}



%%%%%%%%%%%  Picture  %%%%%%%%%%
\\[1em]
\centerline{
\includegraphics[width=2in]{1-8-Discont-iv.png}
\qquad \qquad 
\includegraphics[width=2in]{1-8-Discont-v.png} 
}
\vspace{1em}

\noindent
%%%%%%%%%%%%%%%%%%%%%%%%%

%%%%%%%%%%%  Table  %%%%%%%%%%
In fact, we get the following derivatives:
$$\begin{array}{|c||c|c|c|c|c|c|}
\hline &&&&&& \\[-.9em] 
f(x) & \sin(x) & \cos(x) & \tan(x)  & \sec(x) & \csc(x) & \cot(x)
\\[.2em] \hline &&&&&& \\[-.9em]
f'(x) &\cos(x) & -\sin(x) & \sec^2(x) & \tan(x)\sec(x) & -\cot(x)\csc(x) & -\csc^2(x)
\\[.1em] \hline
\end{array}$$
\\[1em]
%%%%%%%%%%%%%%%%%%%%%%%%%%

-- Example: marginal cost, marginal tax rate (next chapter?).
\\


















